\section{Results and Evaluation}
The model was evaluated on the held-out test set (about 15\% of the data) using the saved best model.

\tablestyle
\begin{table}[h]
\centering
\caption{Performance of the proposed model}
\begin{tabularx}{\textwidth}{|l|X|}\hline
\hdrrow \hdr{Metric} & \hdr{Value} \\
\hline
Final test accuracy & \textbf{85--87\%} \\
\hline
Test loss           & \textbf{0.4214} \\
\hline
Training time       & \textbf{10--12 min} \\
\hline
\end{tabularx}
\end{table}

\subsection{Discussion}
Transfer learning allows a model trained on only about two thousand images to reach high accuracy, because the low-level features learned from ImageNet (edges, textures, shapes) transfer well to waste images. Data augmentation and dropout limit over-fitting, and fine-tuning the last 20 layers adapts the high-level features to the four material classes. Classes with similar appearance, such as glass and transparent plastic, or crumpled paper and light metal, are the most likely to be confused.
